% !TEX root =  Lecture5.tex


%%%%%%%%%%%%%%%%%%%%%%%%%%%%%%%%%%%%
% Recap/Agenda 
%%%%%%%%%%%%%%%%%%%%%%%%%%%%%%%%%%%%
\begin{frame}
\frametitle{Module 1: Multiple Regression}
    \begin{itemize}
    \item \hl{Previously:} using a fitted line --- $R^2$, prediction and extrapolation, and a first look at indicator variables.
    \item \hl{This time:} multiple regression.
      \begin{itemize}
      \item the multiple regression model and conditional interpretation
      \item categorical and indicator variables in a multiple regression model
      \item comparing model fit using $R^2$ and adjusted $R^2$
      \end{itemize}
    \item \hl{Reading:} IMS Chapter 8.1--8.3.
    \item \hl{Homework for today:}
    \item \hl{Homework for next class:}
    \end{itemize}
\end{frame}


%%%%%%%%%%%%%%%%%%%%%%%%%%%%%%%%%%%%
% Learning objectives:
%%%%%%%%%%%%%%%%%%%%%%%%%%%%%%%%%%%%
\begin{frame}
    \frametitle{Lecture Objectives}
    \goals{%
    \begin{itemize}
    \item explain what changes when a model has more than one predictor;
    \item interpret each coefficient while holding the other predictors fixed;
    \item identify reference levels and indicator variables in software output;
    \item explain why adjusted $R^2$, not $R^2$, is used to compare models with different numbers of predictors.
    \end{itemize}}
    \objectives{%
    \begin{itemize}
    \item \textbf{(M1, LO3)} Using Regression Models
    \item \textbf{(M1, LO4)} Model Selection and $R^{2}$ and Adjusted $R^{2}$
    \item \textbf{(M1, LO5)} Categorical Variables and Indicator Variables
    \end{itemize}}
\end{frame}

    